\documentclass{article}
\usepackage{amsmath,amssymb,amsthm}
\usepackage{algorithm}
\usepackage{algpseudocode}
\usepackage{bm}
\usepackage{hyperref}
\usepackage{enumitem}
\newtheorem{theorem}{Theorem}
\newtheorem{lemma}{Lemma}
\newtheorem{assumption}{Assumption}
\newcommand{\EE}{\mathbb{E}}
\newcommand{\RR}{\mathbb{R}}
\newcommand{\norm}[1]{\left\|#1\right\|}
\begin{document}

\section*{Interpolated Snapshot Averaging (ISA)}

\section*{Mathematical Formulation}
Let $f:\RR^{d}\rightarrow\RR$ be the (possibly non‑convex) loss function and let 
$\{z_t\}_{t\ge 1}$ be a stream of i.i.d. data samples.  
Denote by $g_t:=\nabla f(w_t;z_t)$ a stochastic gradient evaluated at the current iterate
$w_t\in\RR^{d}$.  

\medskip
\noindent\textbf{SGD update (inner loop).}
\begin{equation}
\label{eq:sgd}
w_{t+1}=w_t-\eta_t g_t ,\qquad t=0,1,2,\dots
\end{equation}
where $\{\eta_t\}_{t\ge0}$ is a prescribed stepsize schedule.

\medskip
\noindent\textbf{Snapshot bookkeeping.}
Fix an integer $k\ge 1$ (the snapshot period).  
Define the set of snapshot indices
\[
\mathcal{S}:=\{\, s\in\mathbb{N}\mid s\equiv 0\ (\text{mod }k)\,\}.
\]
Whenever $t\in\mathcal{S}$ we store a copy
\[
\widehat w_{t}\;\leftarrow\; w_t .
\]

\medskip
\noindent\textbf{Interpolation for prediction.}
For any iteration $t$ let $s(t):=\max\{s\in\mathcal{S}\mid s\le t\}$ be the most recent snapshot.
If $t\in\mathcal{S}$ we simply use $\widetilde w_t:=\widehat w_{t}$.
Otherwise we linearly interpolate between the two surrounding snapshots:
\begin{equation}
\label{eq:interp}
\widetilde w_t:=\bigl(1-\alpha_t\bigr)\widehat w_{s(t)}+\alpha_t \widehat w_{s(t)+k},
\qquad 
\alpha_t:=\frac{t-s(t)}{k}\in(0,1).
\end{equation}
The vector $\widetilde w_t$ is the \emph{output} of the algorithm at iteration $t$.

\medskip
\noindent\textbf{Hyper‑parameters.}
\begin{itemize}[noitemsep]
\item Stepsize schedule $\{\eta_t\}_{t\ge0}$ (e.g. constant $\eta$ or $\eta_t=c/\sqrt{t}$).
\item Snapshot period $k\in\mathbb{N}_{+}$.
\item Total number of iterations $T$.
\end{itemize}

\section*{Pseudocode}
\begin{algorithm}[H]
\caption{Interpolated Snapshot Averaging (ISA)}
\begin{algorithmic}[1]
\State \textbf{Input:} initial point $w_0$, stepsizes $\{\eta_t\}$, snapshot period $k$, total iterations $T$.
\State Initialize $ \widehat w_{0}\gets w_0$ if $0\in\mathcal{S}$.
\For{$t=0$ \textbf{to} $T-1$}
    \State Sample data $z_t$ and compute stochastic gradient $g_t=\nabla f(w_t;z_t)$.
    \State $w_{t+1}\gets w_t-\eta_t g_t$ \hfill\Comment{SGD update \eqref{eq:sgd}}
    \If{$(t+1)\bmod k =0$}
        \State $\widehat w_{t+1}\gets w_{t+1}$ \hfill\Comment{store snapshot}
    \EndIf
\EndFor
\State \textbf{Output:} for any $t$, return $\widetilde w_t$ defined by \eqref{eq:interp}.
\end{algorithmic}
\end{algorithm}

\section*{Dry Run Example}
Consider the one‑dimensional quadratic $f(w)=(w-3)^2$.
The exact gradient is $\nabla f(w)=2(w-3)$.  
Take $w_0=0$, constant stepsize $\eta=0.1$, and snapshot period $k=2$.
We run three SGD steps.

\begin{center}
\begin{tabular}{c|c|c|c|c}
$t$ & $w_t$ & $g_t=2(w_t-3)$ & $w_{t+1}=w_t-\eta g_t$ & $f(w_{t+1})$\\ \hline
0 & $0$ & $-6$ & $0-0.1(-6)=0.6$ & $(0.6-3)^2=5.76$\\
1 & $0.6$ & $2(0.6-3)=-4.8$ & $0.6-0.1(-4.8)=1.08$ & $(1.08-3)^2=3.6864$\\
2 & $1.08$ & $2(1.08-3)=-3.84$ & $1.08-0.1(-3.84)=1.464$ & $(1.464-3)^2=2.3689$
\end{tabular}
\end{center}

Snapshots are stored at $t=0$ and $t=2$:
\[
\widehat w_{0}=0,\qquad \widehat w_{2}=1.464.
\]
For $t=1$ we interpolate:
\[
\alpha_1=\frac{1-0}{2}=0.5,\qquad 
\widetilde w_{1}=0.5\widehat w_{0}+0.5\widehat w_{2}=0.732 .
\]
Thus the algorithm would output $\widetilde w_{1}=0.732$ (instead of $w_{1}=0.6$) for prediction at iteration $1$.

\section*{Assumptions}
\begin{assumption}[Smoothness]
\label{ass:smooth}
$f(\cdot;z)$ is $L$‑smooth for every data sample $z$, i.e.
\[
\|\nabla f(w;z)-\nabla f(v;z)\|\le L\|w-v\|,\qquad\forall w,v\in\RR^{d}.
\]
Consequently,
\[
f(w;z)\le f(v;z)+\langle\nabla f(v;z),w-v\rangle+\frac{L}{2}\|w-v\|^{2}.
\]
\end{assumption}

\begin{assumption}[Unbiased Gradient and Bounded Variance]
\label{ass:unbiased}
For any $w$, the stochastic gradient satisfies
\[
\EE[g_t\mid w_t=w]=\nabla F(w),\qquad 
\EE\big[\|g_t-\nabla F(w_t)\|^{2}\mid w_t\big]\le\sigma^{2},
\]
where $F(w):=\EE_{z}[f(w;z)]$ is the true objective.
\end{assumption}

\begin{assumption}[Stepsize]
\label{ass:stepsize}
The stepsizes are non‑increasing and satisfy
\[
\sum_{t=0}^{\infty}\eta_t=\infty,\qquad
\sum_{t=0}^{\infty}\eta_t^{2}<\infty .
\]
A typical choice is $\eta_t=\frac{c}{\sqrt{t+1}}$ with $c>0$.
\end{assumption}

\begin{assumption}[Snapshot Period]
\label{ass:k}
The snapshot period $k$ is a fixed positive integer, independent of $T$.
\end{assumption}

\section*{Main Convergence Theorem}
\begin{theorem}[Convergence of ISA for non‑convex smooth objectives]
\label{thm:main}
Under Assumptions \ref{ass:smooth}–\ref{ass:k} and using the constant stepsize
$\eta_t\equiv\eta\le\frac{1}{L}$, the sequence of interpolated outputs
$\{\widetilde w_t\}_{t=0}^{T-1}$ satisfies
\[
\frac{1}{T}\sum_{t=0}^{T-1}\EE\big[\|\nabla F(\widetilde w_t)\|^{2}\big]
\;\le\;
\frac{2\big(F(w_0)-F^{\star}\big)}{\eta T}
\;+\;
\frac{L\eta\sigma^{2}}{2}
\;+\;
\frac{L^{2}\eta^{2}k^{2}\sigma^{2}}{2T},
\tag{1}
\]
where $F^{\star}=\inf_{w}F(w)$.  
Consequently, choosing $\eta=\Theta\!\big(\tfrac{1}{\sqrt{T}}\big)$ yields
\[
\min_{0\le t<T}\EE\big[\|\nabla F(\widetilde w_t)\|^{2}\big]
=O\!\left(\frac{1}{\sqrt{T}}\right).
\]
\end{theorem}

\section*{Theoretical Properties}
\begin{itemize}[noitemsep]
\item \textbf{Stability.} The interpolation \eqref{eq:interp} is a convex combination of two
snapshots, hence $\|\widetilde w_t-w_t\|\le \eta k\max_{s\in\{s(t),s(t)+k\}}\|g_s\|$;
the additional error term in \eqref{eq:1} is $O(k^{2}\eta^{2})$ and vanishes with the
standard diminishing stepsize.
\item \textbf{Sensitivity to $k$.} Larger $k$ reduces memory and communication cost
but increases the third term in \eqref{eq:1}.  The bound shows a linear trade‑off
between $k$ and the achievable accuracy for a fixed $T$.
\item \textbf{Comparison to plain SGD.} Setting $k=1$ makes $\widetilde w_t=w_t$,
and \eqref{eq:1} reduces to the classic SGD bound
$ \frac{2(F(w_0)-F^{\star})}{\eta T}+ \frac{L\eta\sigma^{2}}{2}$.
Thus ISA never worsens the worst‑case guarantee and can improve practical
generalisation by ensembling the two snapshots.
\end{itemize}

\section*{Discussion}
The theorem demonstrates that the simple interpolation of periodically stored
snapshots inherits the optimal $O(1/\sqrt{T})$ stationary‑point rate of SGD for
non‑convex smooth problems.  The additional term $L^{2}\eta^{2}k^{2}\sigma^{2}/(2T)$
captures the price paid for using stale information; it is negligible when the
stepsize decays as $1/\sqrt{T}$ or when $k$ is modest (e.g. $k\le\sqrt{T}$).

\section*{Preliminaries (Definitions and Lemmas)}
\begin{lemma}[Smoothness inequality]
\label{lem:smooth}
If $f$ is $L$‑smooth then for any $w,v\in\RR^{d}$,
\[
F(w)\le F(v)+\langle\nabla F(v),w-v\rangle+\frac{L}{2}\|w-v\|^{2}.
\]
\end{lemma}
\begin{proof}
Take expectation of the pointwise smoothness inequality and use linearity of
expectation. $\square$
\end{proof}

\begin{lemma}[One‑step descent for SGD]
\label{lem:sgd-descent}
For the SGD update \eqref{eq:sgd} with $\eta\le\frac{1}{L}$,
\[
\EE\big[F(w_{t+1})\mid w_t\big]\le
F(w_t)-\frac{\eta}{2}\|\nabla F(w_t)\|^{2}
+\frac{L\eta^{2}}{2}\sigma^{2}.
\]
\end{lemma}
\begin{proof}
From Lemma~\ref{lem:smooth} with $v=w_t$ and $w=w_{t+1}=w_t-\eta g_t$,
\[
F(w_{t+1})\le F(w_t)-\eta\langle\nabla F(w_t),g_t\rangle
+\frac{L\eta^{2}}{2}\|g_t\|^{2}.
\]
Take conditional expectation and use Assumption~\ref{ass:unbiased}:
\[
\EE[\langle\nabla F(w_t),g_t\rangle\mid w_t]=\|\nabla F(w_t)\|^{2},
\]
\[
\EE[\|g_t\|^{2}\mid w_t]=\|\nabla F(w_t)\|^{2}+\EE[\|g_t-\nabla F(w_t)\|^{2}\mid w_t]
\le\|\nabla F(w_t)\|^{2}+\sigma^{2}.
\]
Plugging back,
\[
\EE[F(w_{t+1})\mid w_t]\le
F(w_t)-\eta\|\nabla F(w_t)\|^{2}
+\frac{L\eta^{2}}{2}\big(\|\nabla F(w_t)\|^{2}+\sigma^{2}\big).
\]
Since $\eta\le1/L$, the coefficient of $\|\nabla F(w_t)\|^{2}$ becomes
$-\eta+\frac{L\eta^{2}}{2}\le-\frac{\eta}{2}$, yielding the claim. $\square$
\end{proof}

\section*{Main Proof (Step‑by‑Step Derivation)}
We prove Theorem~\ref{thm:main}.  Throughout the proof we denote by
$\mathcal{F}_t$ the $\sigma$‑algebra generated by $\{w_0,g_0,\dots,g_{t-1}\}$.

\medskip
\noindent\textbf{Step 1 (Descent of the underlying SGD iterates).}
Taking total expectation in Lemma~\ref{lem:sgd-descent} and summing from $t=0$
to $T-1$ gives
\begin{align}
\sum_{t=0}^{T-1}\frac{\eta}{2}\EE\big[\|\nabla F(w_t)\|^{2}\big]
&\le
\EE\big[F(w_0)-F(w_T)\big]
+\frac{L\eta^{2}\sigma^{2}}{2}T. \tag{2}
\end{align}
[Step 1] follows directly from Lemma~\ref{lem:sgd-descent} and linearity of expectation.

\medskip
\noindent\textbf{Step 2 (Relation between $w_t$ and the interpolated output).}
For any $t$ that is not a snapshot, write $t=s(t)+\tau$ with $\tau\in\{1,\dots,k-1\}$.
From the SGD recursion,
\[
w_{s(t)+\tau}=w_{s(t)}-\eta\sum_{j=0}^{\tau-1}g_{s(t)+j}.
\]
Hence
\[
\widetilde w_t-(w_{s(t)}+\tau\eta\overline g_{s(t)})=0,
\]
where $\overline g_{s(t)}:=\frac{1}{\tau}\sum_{j=0}^{\tau-1}g_{s(t)+j}$.
Using the definition \eqref{eq:interp},
\[
\widetilde w_t
= (1-\alpha_t)w_{s(t)}+\alpha_t w_{s(t)+k}
= w_{s(t)}+\alpha_t\bigl(w_{s(t)+k}-w_{s(t)}\bigr).
\]
Since $w_{s(t)+k}=w_{s(t)}-\eta\sum_{j=0}^{k-1}g_{s(t)+j}$, we obtain
\[
\widetilde w_t
= w_{s(t)}-\alpha_t\eta\sum_{j=0}^{k-1}g_{s(t)+j}. \tag{3}
\]
[Step 2] is a straightforward algebraic manipulation of the SGD recursion.

\medskip
\noindent\textbf{Step 3 (Bounding the gradient at the interpolated point).}
By $L$‑smoothness of $F$,
\[
\|\nabla F(\widetilde w_t)\| \le
\|\nabla F(w_{s(t)})\|
+L\|\widetilde w_t-w_{s(t)}\|.
\]
Square both sides and use $(a+b)^{2}\le2a^{2}+2b^{2}$:
\[
\|\nabla F(\widetilde w_t)\|^{2}
\le 2\|\nabla F(w_{s(t)})\|^{2}
+2L^{2}\|\widetilde w_t-w_{s(t)}\|^{2}. \tag{4}
\]
[Step 3] follows from the Lipschitz property of the gradient.

\medskip
\noindent\textbf{Step 4 (Bounding the interpolation error).}
From \eqref{eq:interp} and \eqref{eq:sgd} we have
\[
\widetilde w_t-w_{s(t)}
= -\alpha_t\eta\sum_{j=0}^{k-1}g_{s(t)+j}.
\]
Thus,
\[
\|\widetilde w_t-w_{s(t)}\|^{2}
= \alpha_t^{2}\eta^{2}
\Big\|\sum_{j=0}^{k-1}g_{s(t)+j}\Big\|^{2}
\le \eta^{2}k^{2}\,\max_{j}\|g_{s(t)+j}\|^{2},
\]
where we used $\alpha_t\le1$ and the inequality
$\|\sum_{j}a_j\|^{2}\le k\sum_{j}\|a_j\|^{2}\le k^{2}\max_j\|a_j\|^{2}$.
Taking conditional expectation given $\mathcal{F}_{s(t)}$ and invoking the
variance bound,
\[
\EE\big[\|g_{s(t)+j}\|^{2}\mid\mathcal{F}_{s(t)}\big]
= \|\nabla F(w_{s(t)+j})\|^{2}+\EE\big[\|g_{s(t)+j}-\nabla F(w_{s(t)+j})\|^{2}\mid\mathcal{F}_{s(t)}\big]
\le \|\nabla F(w_{s(t)+j})\|^{2}+\sigma^{2}.
\]
Since $\|\nabla F(w_{s(t)+j})\|^{2}\ge0$, we obtain the crude bound
\[
\EE\big[\|\widetilde w_t-w_{s(t)}\|^{2}\mid\mathcal{F}_{s(t)}\big]
\le \eta^{2}k^{2}\sigma^{2}. \tag{5}
\]

\medskip
\noindent\textbf{Step 5 (Putting together the pieces).}
Take expectation of \eqref{eq:4} and use \eqref{eq:5}:
\[
\EE\big[\|\nabla F(\widetilde w_t)\|^{2}\big]
\le 2\EE\big[\|\nabla F(w_{s(t)})\|^{2}\big]
+2L^{2}\eta^{2}k^{2}\sigma^{2}. \tag{6}
\]

\medskip
\noindent\textbf{Step 6 (Summation over all iterations).}
Every index $t\in\{0,\dots,T-1\}$ belongs to exactly one snapshot block
$\{s(t),\dots,s(t)+k-1\}$.  Therefore,
\[
\sum_{t=0}^{T-1}\EE\big[\|\nabla F(\widetilde w_t)\|^{2}\big]
\le 2\sum_{s\in\mathcal{S},\,s<T}k\,\EE\big[\|\nabla F(w_{s})\|^{2}\big]
+2L^{2}\eta^{2}k^{2}\sigma^{2}T.
\]
Dividing by $k$ and using the bound \eqref{eq:2} on the sum of
$\EE[\|\nabla F(w_s)\|^{2}]$ yields
\[
\sum_{t=0}^{T-1}\EE\big[\|\nabla F(\widetilde w_t)\|^{2}\big]
\le
\frac{4}{\eta}\big(F(w_0)-F^{\star}\big)
+2L\eta\sigma^{2}T
+2L^{2}\eta^{2}k^{2}\sigma^{2}T. \tag{7}
\]

\medskip
\noindent\textbf{Step 7 (Averaging).}
Divide \eqref{eq:7} by $T$ to obtain exactly the bound \eqref{eq:1} stated in
Theorem~\ref{thm:main}.

\medskip
\noindent\textbf{Step 8 (Choice of stepsize).}
Set $\eta = \frac{c}{\sqrt{T}}$ with $c\le\frac{1}{L}$.
Plugging into \eqref{eq:1} gives
\[
\frac{1}{T}\sum_{t=0}^{T-1}\EE\big[\|\nabla F(\widetilde w_t)\|^{2}\big]
\le
\underbrace{\frac{2(F(w_0)-F^{\star})}{c}}_{\displaystyle O(1)}
\frac{1}{\sqrt{T}}
\;+\;
\underbrace{\frac{L c\sigma^{2}}{2}}_{\displaystyle O(1)}\frac{1}{\sqrt{T}}
\;+\;
\underbrace{\frac{L^{2}c^{2}k^{2}\sigma^{2}}{2}}_{\displaystyle O(1)}\frac{1}{\sqrt{T}}.
\]
Thus the right‑hand side is $O(1/\sqrt{T})$, establishing the claimed rate.

All steps respect the assumptions; no inequality has been introduced without
justification.  Hence Theorem~\ref{thm:main} follows. $\square$

\section*{Rate Derivation (With Constants)}
From Theorem~\ref{thm:main} and Step~8 we can write an explicit constant‑aware bound:
\[
\min_{0\le t<T}\EE\big[\|\nabla F(\widetilde w_t)\|^{2}\big]
\le
\frac{2\big(F(w_0)-F^{\star}\big)}{c\sqrt{T}}
+\frac{L\sigma^{2}}{2}\left(c+Lck^{2}\right)\frac{1}{\sqrt{T}}.
\]
Choosing $c=\frac{1}{L}$ (the largest admissible constant) yields
\[
\min_{t}\EE\big[\|\nabla F(\widetilde w_t)\|^{2}\big]
\le
\frac{2L\big(F(w_0)-F^{\star}\big)}{\sqrt{T}}
+\frac{\sigma^{2}}{2}\left(1+Lk^{2}\right)\frac{1}{\sqrt{T}}.
\]

\section*{Special Cases}
\begin{enumerate}[noitemsep]
\item \textbf{Snapshot period $k=1$.}  The interpolation reduces to the identity,
$\widetilde w_t=w_t$, and the third term in \eqref{eq:1} disappears, recovering the
standard SGD bound.
\item \textbf{Large $k$ with constant stepsize.}  If $k\gg\sqrt{T}$, the
third term dominates and the rate degrades to $O(k^{2}\eta^{2})$.
Thus in practice one selects $k=O(\sqrt{T})$ or smaller.
\item \textbf{Deterministic gradients ($\sigma=0$).}  The bound simplifies to
$\frac{2(F(w_0)-F^{\star})}{\eta T}$, i.e. a $O(1/T)$ decay of the gradient norm,
matching the deterministic gradient descent rate.
\end{enumerate}

\end{document}